\documentclass[border=6pt]{standalone}
\usepackage{amsmath}
\usepackage{amssymb}
\usepackage{array}
\usepackage{booktabs}
\begin{document}
\begin{tabular}{@{} >{\bfseries\raggedright\arraybackslash}p{1.6cm} >{\centering\arraybackslash}p{3.4cm} >{\centering\arraybackslash}p{3.4cm} >{\centering\arraybackslash}p{3.4cm} >{\centering\arraybackslash}p{3.2cm} @{}}
\toprule
\multicolumn{5}{c}{\textbf{CONCEPT SUMMARY: Properties of Logarithms}} \\
\midrule
 & \textbf{Product Property} & \textbf{Quotient Property} & \textbf{Power Property} & \textbf{Change of Base} \\
\midrule
ALGEBRA &
$\log_b(mn)=\log_b m+\log_b n$ &
$\log_b\left(\dfrac{m}{n}\right)=\log_b m-\log_b n$ &
$\log_b(m^{n})=n\cdot\log_b m$ &
$\log_b m=\dfrac{\log_a m}{\log_a b}$ \\
\midrule
WORDS &
The log of a product is the sum of the logs. &
The log of a quotient is the difference of the logs. &
The log of a number raised to a power is the power multiplied by the log of the number. &
The log base $b$ of a number is equal to the log base $a$ of the number divided by the log base $a$ of $b$. \\
\midrule
NUMBERS &
$\log_2(20)=\log_2(4)+\log_2(5)$ &
$\log_{10}\left(\dfrac{2}{3}\right)=\log_{10}2-\log_{10}3$ &
$\log_3(16)=4\cdot\log_3 2$ &
$\log_5 7=\dfrac{\log 7}{\log 5}$ \\
\bottomrule
\end{tabular}
\end{document}
